\documentclass[12pt,a4paper]{article}
\usepackage[utf8]{inputenc}
\usepackage[T1]{fontenc}
\usepackage{amsmath,amssymb,amsthm}
\usepackage{booktabs}
\usepackage{geometry}
\usepackage{enumitem}
\usepackage{hyperref}
\usepackage{fancyhdr}
\usepackage{titlesec}

\geometry{margin=1in}
\pagestyle{fancy}
\fancyhf{}
\rhead{Lecture 10}
\lhead{Quantum Systems via Linear Algebra}
\cfoot{\thepage}

\titleformat{\section}{\large\bfseries}{}{0em}{}
\titleformat{\subsection}{\normalsize\bfseries}{}{0em}{}

\newtheorem{definition}{Definition}
\newtheorem{theorem}{Theorem}

\title{\textbf{Lecture 10: Uncertainty, Ehrenfest, and the Classical Limit} \\ \large Fourier Duality, Particles in Potentials, and Engineering Bridges}
\author{Quantum Systems via Linear Algebra, Differential Equations, and Probability}
\date{}

\begin{document}
\maketitle
\thispagestyle{fancy}

% ============================================================
\section{1. The Uncertainty Principle as Fourier Duality}
% ============================================================

\subsection{1.1 Origin}
Position $\hat{x}$ and momentum $\hat{p}$ do not commute (Lecture~9):
\[
    [\hat{x},\hat{p}] = i\hbar.
\]
Therefore no state can be a simultaneous eigenstate of both.

\subsection{1.2 Fourier duality (the Gabor limit)}
\begin{itemize}[nosep]
    \item A position eigenstate $\delta(x-x_0)$ is perfectly localized in $x$ but has equal amplitude for \emph{all} momenta---completely delocalized in $p$.
    \item A momentum eigenstate $e^{ipx}$ is a plane wave spread over all space---completely delocalized in $x$.
    \item In general: the \textbf{narrower} $v(x)$ is in position, the \textbf{broader} $\tilde{v}(p)$ is in momentum, and vice versa.  This is a fundamental property of Fourier transforms.
\end{itemize}

\medskip
\noindent\textbf{Engineering connection:} This is the \textbf{Gabor limit} (bandwidth--duration product).  In signal processing, a signal cannot be simultaneously narrow in time and narrow in frequency: $\Delta t \cdot \Delta f \geq 1/(4\pi)$.  The quantum uncertainty principle $\Delta x\,\Delta p \geq \hbar/2$ is the \emph{same} mathematical inequality applied to position/momentum instead of time/frequency.

\subsection{1.3 Defining uncertainty}
The \textbf{uncertainty} (standard deviation) in position is
\[
    (\Delta x)^2 = \langle \hat{x}^2 \rangle - \langle \hat{x} \rangle^2 = \int |v(x)|^2\,x^2\,dx - \left(\int |v(x)|^2\,x\,dx\right)^2.
\]
Similarly for momentum:
\[
    (\Delta p)^2 = \langle \hat{p}^2 \rangle - \langle \hat{p} \rangle^2 = \int |\tilde{v}(p)|^2\,p^2\,dp - \left(\int |\tilde{v}(p)|^2\,p\,dp\right)^2.
\]

\subsection{1.4 The inequality}

\begin{theorem}[Heisenberg uncertainty principle / Gabor limit]
For any state $v \in L^2(\mathbb{R})$,
\[
    \boxed{\Delta x\;\Delta p \;\geq\; \frac{\hbar}{2}.}
\]
\end{theorem}

\subsection{1.5 Proof}
Assume $\langle \hat{x} \rangle = \langle \hat{p} \rangle = 0$ (shift coordinates if needed).  Define the one-parameter family
\[
    \mathbf{w}(\alpha) = \bigl(\hat{x} + i\alpha\,\hat{p}\bigr)\mathbf{v}, \qquad \alpha\in\mathbb{R}.
\]
The norm-squared is non-negative:
\[
    \|\mathbf{w}(\alpha)\|^2 = \langle \hat{x}^2 \rangle + \alpha\langle i[\hat{x},\hat{p}] \rangle + \alpha^2\langle \hat{p}^2 \rangle \geq 0.
\]
Substituting $[\hat{x},\hat{p}]=i\hbar$:
\[
    (\Delta x)^2 - \alpha\hbar + \alpha^2(\Delta p)^2 \geq 0 \qquad \forall\;\alpha.
\]
This is a quadratic in $\alpha$ that is always non-negative, so its discriminant must be non-positive:
\[
    \hbar^2 - 4(\Delta x)^2(\Delta p)^2 \leq 0 \quad\Longrightarrow\quad
    \Delta x\;\Delta p \geq \frac{\hbar}{2}.\qquad\qed
\]

% ============================================================
\section{2. The Schr\"{o}dinger Equation for a Particle}
% ============================================================

\subsection{2.1 Free particle}
With system matrix $H = \frac{\hat{p}^2}{2m}$ and $\hat{p}=-i\hbar\frac{d}{dx}$ in position space:
\[
    \boxed{i\hbar\frac{\partial v}{\partial t} = -\frac{\hbar^2}{2m}\frac{\partial^2 v}{\partial x^2}.}
\]

\subsection{2.2 Particle in a potential}
Adding a potential $V(x)$:
\[
    H = \frac{\hat{p}^2}{2m} + V(\hat{x}),
\]
the Schr\"{o}dinger equation is
\[
    \boxed{i\hbar\frac{\partial v}{\partial t} = -\frac{\hbar^2}{2m}\frac{\partial^2 v}{\partial x^2} + V(x)\,v(x,t).}
\]

\subsection{2.3 Energy eigenstates}
Separation of variables $v(x,t)=\phi(x)\,e^{-iEt/\hbar}$ yields the time-independent equation:
\[
    -\frac{\hbar^2}{2m}\frac{d^2\phi}{dx^2} + V(x)\,\phi(x) = E\,\phi(x).
\]
The allowed energies $E$ and corresponding eigenfunctions $\phi(x)$ depend on the specific potential $V(x)$.

% ============================================================
\section{3. Ehrenfest's Theorem: The Classical Limit}
% ============================================================

Using the Schr\"odinger equation, one can derive how expectation values evolve:

\subsection{3.1 Velocity equation}
\[
    \boxed{\frac{d}{dt}\langle \hat{x} \rangle = \frac{\langle \hat{p} \rangle}{m}.}
\]
This is exact.  The velocity of the center of the wave packet equals the average momentum divided by mass.

\subsection{3.2 Force equation}
\[
    \boxed{\frac{d}{dt}\langle \hat{p} \rangle = -\left\langle \frac{dV}{dx} \right\rangle.}
\]
This is Newton's second law in expectation-value form.

\subsection{3.3 When does classical mechanics emerge?}
The expectation value $\langle f(\hat{x}) \rangle$ is \emph{not} the same as $f(\langle \hat{x} \rangle)$.  They agree when the wave packet is well-localized compared to the scale on which $V(x)$ varies.  This requires:
\begin{itemize}[nosep]
    \item \textbf{Heavy particles}: large $m$ implies small $\Delta x$.
    \item \textbf{Smooth potentials}: features in $V(x)$ much broader than $\Delta x$.
\end{itemize}

\medskip
\noindent\textbf{Engineering analog:} This is the ray-optics limit---when the wavelength is much smaller than the feature size, waves propagate like rays (particles).

% ============================================================
\section{4. Engineering Bridge Table}
% ============================================================

\medskip
\noindent The following table summarizes the correspondences between quantum mechanics (as presented in these notes) and familiar engineering concepts.

\medskip
\begin{center}
\renewcommand{\arraystretch}{1.3}
\begin{tabular}{@{}p{0.38\linewidth}p{0.55\linewidth}@{}}
\toprule
\textbf{Quantum concept} & \textbf{Engineering analog} \\
\midrule
State vector $\mathbf{v} \in \mathbb{C}^N$ or $L^2(\mathbb{R})$ & Signal / state vector in a state-space model \\
Self-adjoint matrix $M = M^\dagger$ & Symmetric (Hermitian) transfer matrix \\
Eigendecomposition of $M$ & Normal-mode analysis \\
Schr\"odinger ODE $i\hbar\dot{\mathbf{v}} = H\mathbf{v}$ & LTI system $\dot{\mathbf{x}} = A\mathbf{x}$ with $A = -iH/\hbar$ \\
Unitary evolution $U^\dagger U = I$ & Lossless network / power-preserving scattering matrix $S^\dagger S = I$ \\
Energy eigenstates & Normal modes / poles of the system \\
Fourier transform (position $\leftrightarrow$ momentum) & Time-domain $\leftrightarrow$ frequency-domain \\
$\Delta x\,\Delta p \geq \hbar/2$ & Gabor limit $\Delta t\,\Delta f \geq 1/(4\pi)$ \\
Kronecker product $\otimes$ & \texttt{numpy.kron} / composite state space \\
Schmidt decomposition & SVD of the coefficient matrix \\
Partial trace (density matrix) & Marginalization over latent variables \\
Commutator $[A,B] \neq 0$ & Non-commuting transfer functions \\
Born rule $P = |\mathbf{e}_i^\dagger\mathbf{v}|^2$ & Squared projection onto a basis vector \\
\bottomrule
\end{tabular}
\end{center}

% ============================================================
\section{5. Course Summary}
% ============================================================

Over ten lectures, the foundations of quantum mechanics have been developed using linear algebra, ODEs, and probability:

\begin{enumerate}[nosep]
    \item \textbf{States} are unit-norm column vectors in a complex vector space (Axiom~1).
    \item \textbf{Observables} are self-adjoint matrices; their eigenvalues are the possible measurement outcomes (Axiom~2).
    \item \textbf{Born's rule} gives probabilities as $|\mathbf{e}_i^\dagger\mathbf{v}|^2$ (Axiom~3).
    \item \textbf{Time evolution} is unitary, governed by $i\hbar\dot{\mathbf{v}} = H\mathbf{v}$ (Axiom~4).
    \item \textbf{Composite systems} use Kronecker products; entanglement $\Leftrightarrow$ $\operatorname{rank}(\Psi) \geq 2$.
    \item \textbf{The density matrix} describes subsystems via partial trace (marginalization); locality is guaranteed by unitarity.
    \item \textbf{Continuous systems} use $L^2(\mathbb{R})$; the Fourier transform connects position and momentum representations.
    \item The \textbf{uncertainty principle} $\Delta x\,\Delta p\geq\hbar/2$ is Fourier duality (the Gabor limit).
    \item \textbf{Ehrenfest's theorem} recovers Newton's laws for expectation values; classical mechanics emerges when wave packets are narrow compared to potential features.
\end{enumerate}

% ============================================================
\newpage
\appendix
\section*{Appendix: Dirac Notation Translation}
% ============================================================

Physics texts use Dirac's bra-ket notation.  The following table translates between the notation used in these notes and Dirac notation.

\medskip
\begin{center}
\renewcommand{\arraystretch}{1.3}
\begin{tabular}{@{}lll@{}}
\toprule
\textbf{Concept} & \textbf{These notes} & \textbf{Dirac notation} \\
\midrule
State (column vector) & $\mathbf{v}$ & $|\psi\rangle$ (``ket'') \\
Conjugate transpose (row vector) & $\mathbf{v}^\dagger$ & $\langle\psi|$ (``bra'') \\
Inner product & $\mathbf{u}^\dagger\mathbf{v}$ & $\langle\phi|\psi\rangle$ (``bracket'') \\
Outer product / projector & $\mathbf{v}\mathbf{v}^\dagger$ & $|\psi\rangle\langle\psi|$ \\
Basis vector & $\mathbf{e}_i$ & $|i\rangle$ or $|\lambda_i\rangle$ \\
Matrix element & $\mathbf{e}_i^\dagger M\mathbf{e}_j$ & $\langle i|M|j\rangle$ \\
Expectation value & $\mathbf{v}^\dagger M\mathbf{v}$ & $\langle\psi|M|\psi\rangle$ \\
Resolution of identity & $\sum_i \mathbf{e}_i\mathbf{e}_i^\dagger = I$ & $\sum_i |i\rangle\langle i| = \mathbf{1}$ \\
Probability amplitude & $\mathbf{e}_i^\dagger\mathbf{v}$ & $\langle i|\psi\rangle$ \\
Probability & $|\mathbf{e}_i^\dagger\mathbf{v}|^2$ & $|\langle i|\psi\rangle|^2$ \\
Self-adjoint matrix & $M = M^\dagger$ & Hermitian operator $\hat{A} = \hat{A}^\dagger$ \\
Pauli / standard observables & $X, Y, Z$ & $\sigma_x, \sigma_y, \sigma_z$ \\
System matrix (Hamiltonian) & $H$ & $\hat{H}$ \\
Kronecker product & $\mathbf{u} \otimes \mathbf{w}$ & $|u\rangle \otimes |w\rangle$ or $|u,w\rangle$ \\
Position-space components & $v(x)$ & $\psi(x) = \langle x|\psi\rangle$ \\
Momentum-space components & $\tilde{v}(p)$ & $\tilde{\psi}(p) = \langle p|\psi\rangle$ \\
\bottomrule
\end{tabular}
\end{center}

\medskip
\noindent\textbf{Convention note.}  Dirac notation makes the conjugate-linearity of the inner product visually obvious: the ``bra'' $\langle\phi|$ is conjugate-linear, the ``ket'' $|\psi\rangle$ is linear.  In column-vector notation, the same information is carried by the dagger: $\mathbf{u}^\dagger$ is the conjugate transpose.  Both notations encode the same mathematics; the choice is a matter of convention.

\end{document}
